\documentclass[10pt,a4paper]{amsart}
\usepackage{amsmath,amssymb,amsthm,graphicx}
% Retained minimal preamble for the Stacks Project chapter sites-cohomology.tex (Cohomology on Sites).
% Verbatim excerpt of the chapter preamble: Xy-pic setup (native preamble lines 14-18) and the native
% theorem declarations with the mathematical macros (native preamble lines 154-217). Copied mechanically;
% nothing added, removed or redefined. The wrapper class and the amsmath/amssymb/amsthm/graphicx base
% packages are supplied by the assembler, so the class/externaldocument part of the native preamble is omitted.
% For commutative diagrams we use Xy-pic
\usepackage[all]{xy}

% We use 2cell for 2-commutative diagrams.
\xyoption{2cell}
% Theorem environments.
%
\theoremstyle{plain}
\newtheorem{theorem}[subsection]{Theorem}
\newtheorem{proposition}[subsection]{Proposition}
\newtheorem{lemma}[subsection]{Lemma}

\theoremstyle{definition}
\newtheorem{definition}[subsection]{Definition}
\newtheorem{example}[subsection]{Example}
\newtheorem{exercise}[subsection]{Exercise}
\newtheorem{situation}[subsection]{Situation}

\theoremstyle{remark}
\newtheorem{remark}[subsection]{Remark}
\newtheorem{remarks}[subsection]{Remarks}

\numberwithin{equation}{subsection}

% Macros
%
\def\lim{\mathop{\mathrm{lim}}\nolimits}
\def\colim{\mathop{\mathrm{colim}}\nolimits}
\def\Spec{\mathop{\mathrm{Spec}}}
\def\Hom{\mathop{\mathrm{Hom}}\nolimits}
\def\Ext{\mathop{\mathrm{Ext}}\nolimits}
\def\SheafHom{\mathop{\mathcal{H}\!\mathit{om}}\nolimits}
\def\SheafExt{\mathop{\mathcal{E}\!\mathit{xt}}\nolimits}
\def\Sch{\mathit{Sch}}
\def\Mor{\mathop{\mathrm{Mor}}\nolimits}
\def\Ob{\mathop{\mathrm{Ob}}\nolimits}
\def\Sh{\mathop{\mathit{Sh}}\nolimits}
\def\NL{\mathop{N\!L}\nolimits}
\def\CH{\mathop{\mathrm{CH}}\nolimits}
\def\proetale{{pro\text{-}\acute{e}tale}}
\def\etale{{\acute{e}tale}}
\def\QCoh{\mathit{QCoh}}
\def\Ker{\mathop{\mathrm{Ker}}}
\def\Im{\mathop{\mathrm{Im}}}
\def\Coker{\mathop{\mathrm{Coker}}}
\def\Coim{\mathop{\mathrm{Coim}}}

% Boxtimes
%
\DeclareMathSymbol{\boxtimes}{\mathbin}{AMSa}{"02}

%
% Macros for moduli stacks/spaces
%
\def\QCohstack{\mathcal{QC}\!\mathit{oh}}
\def\Cohstack{\mathcal{C}\!\mathit{oh}}
\def\Spacesstack{\mathcal{S}\!\mathit{paces}}
\def\Quotfunctor{\mathrm{Quot}}
\def\Hilbfunctor{\mathrm{Hilb}}
\def\Curvesstack{\mathcal{C}\!\mathit{urves}}
\def\Polarizedstack{\mathcal{P}\!\mathit{olarized}}
\def\Complexesstack{\mathcal{C}\!\mathit{omplexes}}
% \Pic is the operator that assigns to X its picard group, usage \Pic(X)
% \Picardstack_{X/B} denotes the Picard stack of X over B
% \Picardfunctor_{X/B} denotes the Picard functor of X over B
\def\Pic{\mathop{\mathrm{Pic}}\nolimits}
\def\Picardstack{\mathcal{P}\!\mathit{ic}}
\def\Picardfunctor{\mathrm{Pic}}
\def\Deformationcategory{\mathcal{D}\!\mathit{ef}}

% --- appended 2026-10-12 (Chief): the project own definitions of the reference/slogan/history
% display environments, copied mechanically from the shared Stacks preamble lines 8-12
% (\usepackage{verbatim} plus the three \newenvironment lines, which the chapter preamble keeps
% verbatim from the project preamble). Nothing added, removed or redefined: these environments
% carry no mathematical content and the project itself typesets them as comments. Without these
% three lines a counted statement whose own text uses them fails with Environment slogan undefined.
\usepackage{verbatim}
\newenvironment{reference}{\comment}{\endcomment}
\newenvironment{slogan}{\comment}{\endcomment}
\newenvironment{history}{\comment}{\endcomment}
\usepackage{cleveref}
\crefname{originaltheorem}{original theorem}{original theorems}
\Crefname{originaltheorem}{Original theorem}{Original theorems}
\makeatletter
\newlabel{injectives-theorem-K-injective-embedding-grothendieck}{{079P}{source}{}{sourcepointer-1}{}}
\newlabel{injectives-theorem-K-injective-embedding-grothendieck@cref}{{[originaltheorem][0][]079P}{[1][0][]source}{}{}{}}
\makeatother
\crefname{originallemma}{original lemma}{original lemmas}
\Crefname{originallemma}{Original lemma}{Original lemmas}
\makeatletter
\newlabel{derived-lemma-derived-adjoint-functors}{{09T5}{source}{}{sourcepointer-2}{}}
\newlabel{derived-lemma-derived-adjoint-functors@cref}{{[originallemma][0][]09T5}{[1][0][]source}{}{}{}}
\makeatother
\crefname{originallemma}{original lemma}{original lemmas}
\Crefname{originallemma}{Original lemma}{Original lemmas}
\makeatletter
\newlabel{categories-lemma-adjoint-fully-faithful}{{07RB}{source}{}{sourcepointer-3}{}}
\newlabel{categories-lemma-adjoint-fully-faithful@cref}{{[originallemma][0][]07RB}{[1][0][]source}{}{}{}}
\makeatother
\title{The derived category of torsion abelian sheaves is the full subcategory of complexes with torsion cohomology}
\author{The Stacks Project Authors}
\date{stacks-project:sites-cohomology.tex:local-snapshot-20261011; GNU FDL 1.2 or later}
\begin{document}
\maketitle
\noindent\textit{Editorial note.} All delivered mathematical text is the author's own, copied from the pinned Stacks Project chapter file sites-cohomology.tex; no statement or proof step was rewritten, reordered or removed.\par\medskip
\noindent\textit{Editorial note.} The proof is the author's complete argument: divisibility, the quasi-isomorphism statement, the K-injectivity and the adjunction identities are all written out; the closing 'Some details omitted' concerns one identification and is accompanied by the author's reference to a more general result that supplies it, and it is preserved verbatim in the counted span.\par\medskip
\noindent\textit{Editorial note.} Every labelled reference occurring inside the counted spans is carried in assembly.external\_references as a non-load-bearing navigation pointer (same-chapter ones at their definition site, cross-chapter ones at the invocation inside the counted span), so no load-bearing mathematics is replaced by a stub. No bibliography citation occurs inside the counted spans.\par\medskip
\noindent\textit{Editorial note.} The checked chapter list for same-name collisions is the set of stacks chapters staged in this repository: chow, cohomology, divisors, etale-cohomology, flat, more-algebra, more-morphisms, morphisms, spaces-more-morphisms, varieties together with this card's chapter sites-cohomology.\par\medskip
\noindent\textit{Editorial note.} The preamble is the chapter's own retained minimal preamble (sources/stacks-sites-cohomology/retained-preamble.tex), a verbatim excerpt of the native preamble: the Xy-pic setup and the native theorem declarations with the mathematical macros.\par\medskip
\noindent\textit{Allowed standard prerequisites (editorial).}\begin{itemize}
\item Torsion abelian sheaves form a Serre subcategory. Same chapter, the torsion abelian sheaves on a site, which are closed under subobjects, quotients and extensions. Supplies the subcategory A whose derived category is described.
\item The derived adjunction of the inclusion of the full subcategory. Derived Categories, Lemma derived-lemma-derived-adjoint-functors, together with the criterion for adjoint functors with invertible unit and counit. Supplies the equivalence criterion applied at the end.
\end{itemize}\medskip
\ifdefined\hypertarget\hypertarget{sourcepointer-1}{}\fi
\noindent\textit{Source reference.} Original theorem 079P: sources/stacks-sites-cohomology/pinned-original-source.tex, source lines 4260--4260. Navigation pointer invoked at proof line 4260; the counted span contains the pointer only and the cited result is not reproduced inside it, so the pointer is registered as a non-load-bearing navigation stub in the established form for Stacks material. No bibliography citation is involved.\par\medskip
\ifdefined\hypertarget\hypertarget{sourcepointer-2}{}\fi
\noindent\textit{Source reference.} Original lemma 09T5: sources/stacks-sites-cohomology/pinned-original-source.tex, source lines 4274--4274. Navigation pointer invoked at proof line 4274; the counted span contains the pointer only and the cited result is not reproduced inside it, so the pointer is registered as a non-load-bearing navigation stub in the established form for Stacks material. No bibliography citation is involved.\par\medskip
\ifdefined\hypertarget\hypertarget{sourcepointer-3}{}\fi
\noindent\textit{Source reference.} Original lemma 07RB: sources/stacks-sites-cohomology/pinned-original-source.tex, source lines 4277--4277. Navigation pointer invoked at proof line 4277; the counted span contains the pointer only and the cited result is not reproduced inside it, so the pointer is registered as a non-load-bearing navigation stub in the established form for Stacks material. No bibliography citation is involved.\par\medskip
\begin{lemma}
\label{lemma-torsion}
Let $\mathcal{C}$ be a site. Let $\mathcal{A} \subset \textit{Ab}(\mathcal{C})$
denote the Serre subcategory consisting of torsion abelian sheaves.
Then the functor $D(\mathcal{A}) \to D_\mathcal{A}(\mathcal{C})$
is an equivalence.
\end{lemma}

\begin{proof}
A key observation is that an injective abelian sheaf $\mathcal{I}$
is divisible. Namely, if $s \in \mathcal{I}(U)$ is a local section,
then we interpret $s$ as a map
$s : j_{U!}\mathbf{Z} \to \mathcal{I}$ and we apply the
defining property of an injective object to
the injective map of sheaves $n : j_{U!}\mathbf{Z} \to j_{U!}\mathbf{Z}$
to see that there exists an $s' \in \mathcal{I}(U)$ with $ns' = s$.

\medskip\noindent
For a sheaf $\mathcal{F}$ denote $\mathcal{F}_{tor}$ its torsion subsheaf.
We claim that if $\mathcal{I}^\bullet$ is a complex of injective abelian
sheaves whose cohomology sheaves are torsion, then
$$
\mathcal{I}^\bullet_{tor} \to \mathcal{I}^\bullet
$$
is a quasi-isomorphism. Namely, by flatness of $\mathbf{Q}$ over $\mathbf{Z}$
we have
$$
H^p(\mathcal{I}^\bullet) \otimes_\mathbf{Z} \mathbf{Q} =
H^p(\mathcal{I}^\bullet \otimes_\mathbf{Z} \mathbf{Q})
$$
which is zero because the cohomology sheaves are torsion.
By divisibility (shown above) we see that
$\mathcal{I}^\bullet \to \mathcal{I}^\bullet \otimes_\mathbf{Z} \mathbf{Q}$
is surjective with kernel $\mathcal{I}^\bullet_{tor}$.
The claim follows from the long exact sequence of cohomology sheaves
associated to the short exact sequence you get.

\medskip\noindent
To prove the lemma we will construct right adjoint
$T : D(\mathcal{C}) \to D(\mathcal{A})$. Namely, given $K$
in $D(\mathcal{C})$ we can represent $K$ by a K-injective complex
$\mathcal{I}^\bullet$ whose cohomology sheaves are injective, see
Injectives, Theorem
\ref{injectives-theorem-K-injective-embedding-grothendieck}.
Then we set $T(K) = \mathcal{I}^\bullet_{tor}$, in other words,
$T$ is the right derived functor of taking torsion.
The functor $T$ is a right adjoint to
$i : D(\mathcal{A}) \to D_\mathcal{A}(\mathcal{C})$.
This readily follows from the observation that
if $\mathcal{F}^\bullet$ is a complex of torsion sheaves, then
$$
\Hom_{K(\mathcal{A})}(\mathcal{F}^\bullet, I^\bullet_{tor}) =
\Hom_{K(\textit{Ab}(\mathcal{C}))}(\mathcal{F}^\bullet, I^\bullet)
$$
in particular $\mathcal{I}^\bullet_{tor}$ is a K-injective complex
of $\mathcal{A}$. Some details omitted; in case of doubt, it also
follows from the more general
Derived Categories, Lemma \ref{derived-lemma-derived-adjoint-functors}.
Our claim above gives that $L = T(i(L))$ for $L$ in $D(\mathcal{A})$
and $i(T(K)) = K$ if $K$ is in $D_\mathcal{A}(\mathcal{C})$.
Using Categories, Lemma \ref{categories-lemma-adjoint-fully-faithful}
the result follows.
\end{proof}

\end{document}
